The following subchapters explain the development work necessary for the implementation of the \acrshort{tof} demonstrator.
In particular, they discuss how the electronics can be constructed. In addition, further
challenges of such a system and how they can be overcome are shown.

\subsection{Signal chain}

Some explanations can be found below on how the \acrshort{tof} measurement is realized.

A block diagram of the overall project, consisting of the electrical and the optical part, is shown in
Figure~\ref{fig:blockdiagram}.

\begin{figure}[H]
    \centering
    \includegraphics[width=0.8\textwidth]{../diagrams/blockdiagram.pdf}
    \caption{Block diagram}\label{fig:blockdiagram}
\end{figure}

A more detailed sketch of the optical signal chain can be seen in the figure~\ref{fig:konzept_signalpfad}.

\begin{figure}[H]
    \centering
    \includegraphics[width=\textwidth]{../diagrams/konzept_signalpfad.pdf}
    \caption{Conceptual signal path}\label{fig:konzept_signalpfad}
\end{figure}

An \acrshort{mcu} takes care of starting and evaluating the measurement results. Generating a start signal
a laser diode is triggered by a laser driver. This pulse, in the form of a light pulse after the laser diode,
is then reflected by a target that is at a variable distance from the electronics.

The reflected light pulse is focused by an optical system and then converted into a photocurrent. This photocurrent is
converted into a voltage by a transimpedance amplifier. Before the \acrshort{mcu} can evaluate this voltage, it must
first be digitized. This is done using a comparator. This can be understood as a 1-bit \acrshort{adc}, since only a
distinction between the two states \dq light\dq\ and \dq no light\dq\ is necessary. The result is a received pulse.

If the time between the start signal and the received pulse is measured, it is possible to draw a conclusion about the variable
distance between the electronics and the target. The time measurement is subject to the highest requirements. As
symbolization: light travels a distance of 30 cm within a nanosecond. The next subchapter explains how such a time measurement
can be realized.

\subsection{Time measurement}

In order for a direct \acrshort{tof} signal to be evaluated, it requires a high-resolution
time measurement. The TDC7200 from Texas Instruments offers a good solution for this.
According to the data sheet \cite{ti2016tdc7200_datasheet}, it can resolve up to 55 ps with a
standard deviation of 35 ps. If we now insert the speed of light and this minimum
temporal resolution in a motion equation, the module should achieve the following spatial resolution according to the formula~\ref{eq:tdc_max_resolution}:

\begin{equation}\label{eq:tdc_max_resolution}
    s_{min} = c \cdot t = 299.8 \cdot 10^6~\frac{m}{s} \cdot 55~ps = 16.5~mm
\end{equation}
\myequations{Maximum spatial resolution of the TDC7200}

In principle, the \acrshort{tof} functionality should be put into operation in two steps.
In the first part, the aim is purely to get to know the TDC7200 \acrshort{ic} better.

Note: In reality, such a resolution will not be achievable. The challenges
consist, for example, in the fact that various delays are present in the entire signal chain.
For example, the switching on of the laser diode, the bandwidth of the amplifier and
the switching time of the comparator play a major role. Also, the jitter of the \acrshort{mcu} when switching on and off
of their outputs will be important at the minimum resolution. In the chapter~\ref{sec:electrical_measurements},
various measurements are used to show how close the TDC7200's specified limit is approached.

\subsection{Approach}\label{sec:approach}

As already explained, it can be assumed that the entire signal path introduces an inaccuracy into the measurement that
is far outside that of the TDC7200. It therefore makes sense to examine this in advance.
For this purpose, a second TDC7200 is used on the electronics, which is responsible for measuring a purely electrical signal
.
\begin{figure}[H]
    \centering
    \includegraphics[width=\textwidth]{../diagrams/konzept_tdc_electrical.pdf}
    \caption{Concept operation \acrshort{tdc} purely electric}\label{fig:konzept_tdc_electrical}
\end{figure}

Figure~\ref{fig:konzept_tdc_electrical} shows a possible design for such a circuit. The wavy
red line symbolizes a variable measurement path. Cables of different lengths can be connected here via a plug.
This makes it possible to verify the measurement accuracy of the TDC7200 purely electrically.

The following points can be checked with such an extension:

\begin{itemize}
    \item Measurement of jitter at the \acrshort{mcu} outputs
    \item Configuration and operation of the TDC7200 (e.g. different measurement modes)
    \item Evaluation of the TDC7200 results
\end{itemize}
